\section{The inherited interface and the new local task}
\label{sec:interface}

\subsection{Precisely which contraction problem is improved}
The root input is a bipartite Boolean tensor network with at most $s$ wires,
maximum vertex degree at most $\lceil s^{9/10}\rceil$, and entries in
\[
 \mathcal B_{P_0,\ell}
   =(\mathbb Z/2^{P_0}\mathbb Z)[V]/(V^\ell),\qquad P_0,\ell\le s.
\]
Entry routines and descriptions satisfy the source interface: polynomial
space and $(s+\mathcal L)^{O(1)}2^{O(s^{9/10})}$ time, with fixed constants.
The contraction routine works in a degree-$O(\log s)$ unramified extension
of the coefficient ring. Descriptions carry immutable creation precisions:
an old entry is always evaluated using its own creation-time representatives,
then lifted to the current precision. Raising precision does not recompute a
previous contraction or silently change a lift~\cite{openai,asymmetric}.

An active call has $M$ wires, a fixed number $r$ of groups, retained valuation
range $u$, cumulative packet exponent $a$, packet length $L=2^a$, and a
retained valuation $i<u$. Write $D=2^{a-i}$. The inputs taken from the reviewed
papers are the following, with their quantitative premises checked for our
schedule in Section~\ref{sec:schedule}:
\begin{enumerate}
\item Nearly lossless asymmetric packet packing gives padded mass
$(1+o(1))M$ and child size at most $(2/r+o(1))M$.
\item Full-block integral Cauchy separation gives, for each directed pair
of different groups, at most $(2+o(1))M/(r^2D)$ families of each mode index
$q\ge0$, with guaranteed weights
\[
 w(i,q)=a+1-i+4q-2\pc(q).
\]
\item Sharp tapered Hermite recovery, with an additive precision guard,
requires local precision $P\le(6+o(1))MD/u$.
\item For every fixed collection of mode exponents, the deformed contraction
factors into group values $Z_g$. A value of $Z_g$ is a contraction of a
specified smaller network, not a free oracle.
\end{enumerate}
The source Cauchy deformation, integrality, recovery, exact group-factor
identity, and packet-entry construction remain mathematical dependencies.
Our proof changes how the resulting finite graphical sum is evaluated and
how the costs are propagated. It is a theorem relative to these explicit
interfaces, not an independent proof of the primary ETH result.

\subsection{Exact budgets and finite index sets}
Allocate a nonnegative integer budget $b_{gh}$ to every directed pair,
with $\sum_{g\ne h}b_{gh}<P$. Let $\mathcal I_{gh}$ be the list of its mode
exponent tuples having \emph{exact} total guaranteed weight $b_{gh}$.
An empty list makes that allocation contribute zero. For all other allocations,
put
\[
 h_{gh}=\log_2|\mathcal I_{gh}|,\qquad
 g_{\{g,h\}}=h_{gh}+h_{hg},\qquad
 \K(M,u)=\kap\frac M{\sqrt u},\quad
 \kap=\frac{\pi\sqrt2}{\ln2}.
\]
The value to compute is the exact finite sum
\begin{equation}
 \sum_{\xi_{gh}\in\mathcal I_{gh}}
       \prod_{g=1}^r Z_g((\xi_e)_{e\ni g}).\label{eq:graphsum}
\end{equation}
Because our $r$ belongs to the fixed set $\{4,7,8,9,12\}$, there are only
polynomially many budget allocations. Exact list cardinalities, enumeration,
and ranking/unranking use the integer knapsack recurrence
\[
 N(j,b)=N(j-1,b)+N(j,b-w_j),\qquad N(0,b)=\mathbf1_{b=0},
\]
with zero values at negative budgets. The bit length of every count is
polynomial in $s$. Choosing a contraction plan can therefore use exact
integer cardinalities; it does not need floating-point logarithms.

\begin{lemma}[Weighted mode entropy]\label{lem:entropy}
For each fixed nonnegative vector $z=(z_{gh})$ on directed pairs,
\begin{equation}
 \sum_{g\ne h}z_{gh}h_{gh}
 \le \left(\frac{\|z\|_2}{r}+o(1)\right)\K(M,u).
 \label{eq:entropy}
\end{equation}
The error is uniform over our fixed finite support families, all exact
budget allocations, and all active depths of the schedule below.
\end{lemma}
\begin{proof}
This is the inherited full-Cauchy generating-function bound~\cite{joint,tree}.
The relevant one-family series obeys
\[
 \sum_{q\ge0}-\ln(1-e^{-\theta(1+4q-2\pc(q))})
  =\frac{\pi^2}{24\theta}+O(\log^2(2/\theta)).
\]
For pair $e$ with $z_e>0$, use parameter $\theta/z_e$ in the exponential
counting bound, multiply the logarithm by $z_e$, and sum. The leading bound is
\[
 \theta P+\frac{\pi^2}{24\theta}
          \frac{(2+o(1))M}{r^2D}\sum_ez_e^2.
\]
Optimizing $\theta$ and using $P\le(6+o(1))MD/u$ gives
$\pi\sqrt2 M\|z\|_2/(r\sqrt u)$ in natural logarithms.
Divide by $\ln2$. Zero weights are omitted. Our guard ensures
$D\ge\lambda^6$ and the saddle has scale $\sqrt u/D$, so the error from
the displayed series is lower order uniformly. This argument does not assume
an independently asymptotic approximation for each small exact list.
\end{proof}

\subsection{Complete operation costs, including table construction}
For a vertex subset $A$, its table sums all internal indices and retains all
indices from $A$ to its complement. Merging disjoint sets $A,B$, with
$O=[r]\setminus(A\cup B)$, is a rectangular matrix product. Its three
logarithmic dimensions are the sums $H_1,H_2,H_3$ of pair entropies between
the classes $A,O$, $A,B$, and $B,O$. With a matrix exponent strictly below
$12/5$, square blocking gives the upper exponent
\begin{equation}
 H_1+H_2+H_3-\frac35\min(H_1,H_2,H_3)
 =\max_j\left(\frac25H_j+\sum_{k\ne j}H_k\right).\label{eq:mm}
\end{equation}
The final contraction of three blocks has the same triangle cost. The trace
or final inner product is no more expensive. The bound pays for intermediate
block tables as well as the final multiplication.

The field-to-ring bridge is the one already proved in the contraction-tree
packet~\cite{tree,tree-old}. A coefficient-ring entry is a vector of
polynomially many integer coordinates. Lift the coordinates, compute their
matrix products modulo enough nonexceptional odd primes, reconstruct exact
integer products by CRT, apply the ring structure constants, and reduce
modulo $2^P$. The required reconstruction bit lengths are polynomial in
$s$, including the logarithm of an exponential matrix dimension. A fixed
rational bilinear matrix algorithm with exponent below $12/5$ supplies
constant-size identities; denominators exclude only finitely many primes.
Published matrix multiplication bounds suffice for this conservative
exponent~\cite{mm}. We do not assume that the original quotient ring is a
field or divide by two or a nilpotent in that ring.

Sequential multiplication and reconstruction use a polynomial factor times
the input/output table storage. Formula~\eqref{eq:mm} is at least the sum
of any two $H_i$, so it bounds all tables in that operation. Sorting records
by their index tuples and scanning them costs only polynomial factors in the
record and index lengths beyond the enumerated table size. This avoids an
unaccounted unit-cost random-access assumption for exponential tables.

\subsection{Operation supports and the prior averaging lemma}
Multiply each of the three linear forms in~\eqref{eq:mm} by five.
The resulting vector $v$, indexed by unordered group pairs, has entries
$0$ inside a class, $2$ on the discounted cross-class, and $5$ on the other
two cross-classes. For a complete tree plan $j$, let $\mathcal V_j$ be its
set of supports and put
\[
 F_j(g)=\frac15\max_{v\in\mathcal V_j}v\cdot g.
\]
A plan first builds three nonempty blocks by balanced binary merges and
then contracts their triangle. It has $r-3$ merges and one triangle,
hence $3(r-2)$ supports. The following is inherited from the latest tree
packet, and is restated to make the new geometry's role explicit.

\begin{lemma}[Support averaging]\label{lem:average}
For $q$ fixed plans, define
\[
 Q=\max_{v_j\in\mathcal V_j}\left\|\sum_{j=1}^qv_j\right\|_2^2.
\]
Choosing the least-cost plan gives nonrecursive time and local storage
bounded, up to polynomial factors, by
\[
 2^{(\beta_r+o(1))\K(M,u)},\qquad
 \beta_r=\frac{\sqrt{2Q}}{5qr}.
\]
\end{lemma}
\begin{proof}
Pointwise,
$\min_jF_j(g)\le q^{-1}\sum_jF_j(g)
=(5q)^{-1}\max_{v_j}(\sum_jv_j)\cdot g$.
Give both orientations of an edge the weight
$(\sum_jv_j)_e/(5q)$. Its squared directed norm is at most
$2Q/(25q^2)$; apply Lemma~\ref{lem:entropy}. The family is fixed and finite.
Every input or output table is bounded by its operation cost as just shown.
This controls storage, including singleton caches, but not the time spent
computing their entries; that time is charged next.
\end{proof}

Each singleton group has $2(r-1)$ incident directed pairs. Therefore its
complete cache requires at most $2^{(\alpha_r+o(1))\K(M,u)}$ recursive
child evaluations, where
\[
 \alpha_r=\frac{\sqrt{2(r-1)}}r.
\]
The fully paid recurrence is
\begin{equation}
 T_d\le\poly(s)\left[
 2^{(\alpha_{r_d}+o(1))K_d}T_{d+1}
 +2^{(\beta_{r_d}+o(1))K_d}\right],\qquad K_d=\K(M_d,u_d).
 \label{eq:paid}
\end{equation}
Summing only the local aggregation costs would miss the first term and
would not establish our result.
